% =============================================================================
%  Irradiation Studies and Design Optimization of the ATLAS Tile Calorimeter
%  for the High-Luminosity LHC
% =============================================================================

\documentclass[a0paper,portrait, blockverticalspace=4em, colspace=2em]{tikzposter}
\tikzposterlatexaffectionproofoff
\usepackage{graphicx}
\usepackage{amsmath}
\usepackage{amssymb}
\usepackage{hyperref}
\usepackage{multicol}
\usetikzlibrary{calc}

\definecolor{maincolor}{HTML}{9e2b2f}
\definecolor{secondarycolor}{RGB}{158, 43, 47}
\definecolor{lightred}{RGB}{237, 199, 201}

\definecolor{tid}{RGB}{254,128,2}   % orange
\definecolor{niel}{RGB}{0,140,0}    % green
\definecolor{see}{RGB}{190,0,62}    % red
\definecolor{sel}{RGB}{0,90,170}    % blue

\newcommand{\looseitems}{\small\begin{itemize}\setlength{\itemsep}{0.3em}}
\newcommand{\tightitems}{\small\begin{itemize}\setlength{\itemsep}{0.15em}}
\newcommand{\captiontext}[1]{{\footnotesize\renewcommand{\baselinestretch}{0.88}\selectfont#1}}

\usetheme{Default}
\definecolorstyle{atlascolors}{
    \colorlet{backgroundcolor}{white}
    \colorlet{titlefgcolor}{white}
    \colorlet{titlebgcolor}{maincolor}
    \colorlet{blocktitlefgcolor}{white}
    \colorlet{blocktitlebgcolor}{maincolor}
    \colorlet{blockbodyfgcolor}{black}
    \colorlet{blockbodybgcolor}{lightred!25}
}{}
\usecolorstyle{atlascolors}

\definetitlestyle{sampletitle}{width=760mm, roundedcorners=20, linewidth=2pt, innersep=10pt, titletotopverticalspace=6mm, titletoblockverticalspace=8mm}{
  \begin{scope}[line width=\titlelinewidth, rounded corners=\titleroundedcorners]
    \draw[color=blocktitlebgcolor, fill=titlebgcolor] (\titleposleft,\titleposbottom) rectangle (\titleposright,\titlepostop);
  \end{scope}
  \node[anchor=east, fill=white, rounded corners=10pt, inner sep=10pt, xshift=-22mm] at ($(\titleposright,\titlepostop)!0.5!(\titleposright,\titleposbottom)$) {
\includegraphics[height=8.6cm]{logos/atlas_transparent.png}};
}
\usetitlestyle{sampletitle}

\title{\parbox{0.75\linewidth}{\textbf{Component-Level Qualification Results\\[2mm]of the ATLAS Tile Calorimeter for the HL-LHC}}}
\author{ATLAS Collaboration}
\institute{Stockholm University, CERN}

\begin{document}
\maketitle[width=0.95\textwidth]

\begin{columns}

% ===========================================================================
\column{0.333}

\block{Introduction \& HL-LHC Upgrade}{
    \looseitems
        \item The ATLAS Tile Calorimeter (TileCal) upgrade is absolutely critical for managing the unprecedented radiation levels and the significantly increased pileup anticipated during the High-Luminosity LHC (HL-LHC) era.
        \item The legacy electronics are being entirely replaced to ensure flawless and continuous operation for the full 10-year lifespan of the upgrade, managing approximately $\sim$35\,Tbps of data across 5376 high-speed optical links.
        \item A central and critical hub of this architecture is the \textbf{Link \& Control Daughterboard (DB6)}. It interfaces up to 12 PMT channels via a complex 400-pin FMC connector directly to the off-detector systems.
    \end{itemize}
}

\block{Daughterboard (DB6) Architecture}{
    \looseitems
        \item System redundancy is heavily prioritized: two functionally-equivalent and independent halves utilize \textbf{Triple Modular Redundancy (TMR)} and Xilinx \textbf{Soft Error Mitigation (SEM)} IP to correct soft errors on the fly.
        \item The DB6 incorporates carefully selected commercial-off-the-shelf (COTS) Kintex UltraScale (KU) and Microsemi ProASIC3 FPGAs, seamlessly integrated with CERN's custom radiation-hard GBTx ASICs.
    \end{itemize}
    \vspace{0.12cm}
    \begin{center}
        \includegraphics[width=1.0\linewidth]{assets/fig_db6_blockdiagram.png}
    \end{center}
    \vspace{-0.2cm}
    \captiontext{Detailed DB6 architecture showing redundancy layers, FPGAs, GBTx ASICs, and optical links.}
}

\block{Front-End Minidrawer Integration}{
    \looseitems
        \item The DB6 interfaces with up to four Minidrawers, handling 12 PMT signals each.
        \item It provides high-precision clock distribution and gathers digital data from the ADCs.
    \end{itemize}
}

\block{Single Event Upset (SEU) Rates}{
    \looseitems
        \item KU FPGAs were irradiated with 54\,MeV protons up to $2.2\times10^{11}$\,p/cm$^2$, exhibiting an SEU rate of 116\,SEU/($10^9$\,p/cm$^2$).
        \item In a worst-case scenario, this yields roughly \textbf{5498 SEUs per year}.
        \item \textbf{Mitigation}: Approximately 5487 of these are correctable by SEM and will not impact nominal runs due to TMR.
        \item The remaining 11 uncorrectable errors per year are managed by TMR and cleared by reloading the FPGA configuration from FLASH during data acquisition stops.
    \end{itemize}
    \vspace{0.12cm}
    \begin{center}
        \includegraphics[width=1.0\linewidth]{assets/fig_diodes_niel_test.pdf}
    \end{center}
    \vspace{-0.2cm}
    \captiontext{Diode forward voltage shifts under NIEL testing.}
}

\block{References}{
    \small
    \begin{enumerate}\setlength{\itemsep}{0.2em}
        \item ATLAS Collaboration, \emph{Technical Design Report for the Phase-II Upgrade of the ATLAS Tile Calorimeter}, CERN-LHCC-2017-019 (2018).
        \item E. Valdes \textit{et al.}, \emph{Radiation studies on the ATLAS Link and Control Daughterboard}, IEEE NSS MIC RTSD (2025).
        \item E. Valdes \textit{et al.}, \emph{Radiation studies performed on the High-Luminosity ATLAS TileCal link Daughterboard}, JINST \textbf{18} (2023) C04011.
    \end{enumerate}
}

% ===========================================================================
\column{0.333}

\block{Simulated Radiation Environment}{
    \looseitems
        \item Radiation qualifications incorporate standard ATLAS collaboration safety factors: 1.5 for simulation uncertainties, and 3.0 to account for lot-to-lot variations (specifically for TID/SEE).
        \item \textbf{\textcolor{tid}{TID Requirement}}: A minimum 108\,Gy threshold is mandated for all active DB6 components.
        \item \textbf{\textcolor{niel}{NIEL Requirement}}: The baseline is set to $13.16\times10^{12}$\,n/cm$^2$, which already incorporates an additional safety factor of 1.3 to compensate for the use of non-monoenergetic testing sources.
        \item \textbf{\textcolor{see}{SEE Requirement}}: A rigorous $4.54\times10^{11}$\,h/cm$^2$ expected fluence limit to ensure negligible data corruption.
    \end{itemize}
}

\block{SFP+ Optical Transceiver Matrix}{
    \tightitems
        \item Reliable, high-bandwidth data transmission at 9.6\,Gbps is absolutely critical for the readout system. 
        \item SFP+ components were evaluated under severe \textcolor{niel}{NIEL} conditions. Out of four tested variants, only two passed the mandatory $14\times10^{12}$\,n/cm$^2$ target.
        \item Additional \textcolor{tid}{TID} and \textcolor{sel}{SEL} testing needs to be carried out on the qualified SFP+ variants before their final integration into the DB6 production design.
    \end{itemize}
    \vspace{0.3cm}
    \begin{center}
    \small
    \begin{tabular}{l l c}
    \hline\hline
    \textbf{SFP+ Model} & \textbf{TID / NIEL Exposure Limit} & \textbf{Status} \\
    \hline
    AVAGO AFBR-709SMZ & Failed during early irradiation & \textcolor{red}{Rejected} \\
    CORETEK CT-A000NPP & 1/9 units failed @ $9\times 10^{12}$ & \textcolor{red}{Rejected} \\
    Direktronik 33-4248 & Passed 21$\times10^{12}$\,n/cm$^2$, $>$200\,Gy & \textcolor{niel}{Qualified} \\
    FS-Europe SFP-10GSR & Passed 21$\times10^{12}$\,n/cm$^2$, $>$200\,Gy & \textcolor{niel}{Qualified} \\
    \hline\hline
    \end{tabular}
    \end{center}
}

\block{Oscillator Breakdown \& Assessment}{
    \tightitems
        \item Timing components were exposed to extended radiation profiles to test long-term frequency stability and phase decoupling effects.
        \item The LTC6902IMS exhibited inconsistent post-irradiation behavior across the tested sample set, leading to its disqualification.
    \end{itemize}
    \vspace{0.3cm}
    \begin{center}
    \small
    \begin{tabular}{p{0.35\linewidth} p{0.3\linewidth} p{0.15\linewidth}}
    \hline\hline
    \textbf{Component} & \textbf{Test Outcome} & \textbf{Status} \\
    \hline
    LTC6902IMS QUAD & 4/15 units failed outright & \textcolor{red}{Removed} \\
    540BAA100M000BBG & Passed all stress tests & \textcolor{niel}{Qualified} \\
    \hline\hline
    \end{tabular}
    \end{center}
}

\block{MOSFET Shifts \& SEL Mitigation}{
    \tightitems
        \item The originally selected FDFMA3N109 MOSFETs displayed notable threshold voltage shifts. These shifts were directly correlated to the \textbf{195\,Gy} Gamma component measured in the reactor cavity.
        \item Consequently, the supplementary over-current protection circuitry—which relied heavily on these MOSFETs for \textcolor{sel}{SEL} mitigation—had to be entirely removed from the final DB6 design iteration.
    \end{itemize}
    \vspace{0.12cm}
    \begin{center}
        \includegraphics[width=1.0\linewidth]{assets/fig_mosfet_threshold.png}
    \end{center}
    \vspace{-0.2cm}
    \captiontext{Threshold voltage shifts depicted in IV curves of irradiated FDFMA3N109 MOSFETs.}
}

\block{DB6 TID Performance}{
    \tightitems
        \item Point-of-load regulator currents and FPGA temperatures remained highly stable over the full $\sim$400\,Gy proton exposure at IFJ Krak\'ow and the 220\,Gy test at CERN CC60.
    \end{itemize}
    \vspace{0.12cm}
    \begin{center}
        \includegraphics[width=1.0\linewidth]{assets/fig_db6_tid_test.png}
    \end{center}
    \vspace{-0.2cm}
    \captiontext{DB6 board currents and KU FPGA temperatures remained stable over the $\sim$400\,Gy \textcolor{tid}{TID} test.}
}

\block{DB6 NIEL Performance}{
    \tightitems
        \item DB6-0 was strongly irradiated at the TRIGA reactor up to the $14\times10^{12}$\,n$_\text{eq}$/cm$^2$ target.
        \item The KU FPGAs showed excellent intrinsic tolerance without significant logic degradation.
        \item The ProASIC and FLASH memory components exhibited minor, temporary reprogrammability after-effects, but recovered completely following standard thermal annealing.
    \end{itemize}
}

% ===========================================================================
\column{0.333}

\block{Regulator Qualification Matrix}{
    \looseitems
        \item Maintaining stable, noise-free power delivery is an absolute priority. Both the LVPS bricks and the DB6 point-of-load (POL) conversion components were rigorously and systematically evaluated across multiple test facilities.
    \end{itemize}
    \vspace{0.3cm}
    \begin{center}
    \small
    \begin{tabular}{p{0.3\linewidth} p{0.35\linewidth} p{0.18\linewidth}}
    \hline\hline
    \textbf{Regulator} & \textbf{Radiation Outcome} & \textbf{Status} \\
    \hline
    LTM4691 DC-DC & Testing scheduled 2025--2026 & \textcolor{tid}{Pending} \\
    LT3080 LDO & Passed combined TID/SEE runs & \textcolor{niel}{Qualified} \\
    DB6 POL Regs. & Output currents perfectly stable & \textcolor{niel}{Qualified} \\
    \hline\hline
    \end{tabular}
    \end{center}
}

\block{LVPS Brick Optimization}{
    \looseitems
        \item The Low-Voltage Power Supply (LVPS) bricks are engineered to operate as efficient 300\,kHz two-switch forward converters, reducing the bulk 200\,V input down to a heavily regulated 10\,V output.
        \item Strategic replacement of the STB57N65M5 MOSFET with the much superior \textbf{SIHFS9N60A} variant yielded a remarkable efficiency increase from $\sim$58\% to \textbf{$\sim$72\%}.
    \end{itemize}
    \vspace{0.12cm}
    \begin{center}
        \includegraphics[width=1.0\linewidth]{assets/fig_lvps_brick.pdf}
    \end{center}
    \vspace{-0.2cm}
    \captiontext{Detailed LVPS brick topology and physical hardware implementation overview.}
    \vspace{0.12cm}
    \begin{center}
        \includegraphics[width=1.0\linewidth]{assets/fig_diodes_tid_test.pdf}
    \end{center}
    \vspace{-0.2cm}
    \captiontext{Diodes TID testing results showing minimal shifts.}
}

\block{Simplified LVPS Wiring Scheme}{
    \looseitems
        \item Long-distance control signals originating from the service cavern were halved from \textbf{18} directly down to \textbf{9} per module.
        \item A novel tri-state logic scheme (on/off/hold) was implemented, significantly easing the complex cable routing and yielding a major reduction in the overall material budget.
    \end{itemize}
    \vspace{0.12cm}
    \begin{center}
        \includegraphics[width=1.0\linewidth]{assets/fig_wire_reduction.pdf}
    \end{center}
    \vspace{-0.2cm}
    \captiontext{Diagram of the 18-to-9 control-wire reduction scheme designed for the LVPS.}
}

\block{LVPS Radiation Campaign}{
    \looseitems
        \item Critical components (LT1681, LTC6241, IR2110) passed TID/NIEL tests with minor shifts. System-level validation occurred at CERN CHARM.
    \end{itemize}
    \vspace{0.12cm}
    \begin{center}
        \includegraphics[width=1.0\linewidth]{assets/fig_lvps_tests_overview.pdf}
    \end{center}
    \vspace{-0.2cm}
    \captiontext{Facilities and beam configurations for the LVPS campaign.}
}

\block{Conclusions \& Mass Production}{
    \looseitems
        \item The extensive and multi-layered testing campaign confirmed full HL-LHC compatibility for the vast majority of DB6 and LVPS electronic components.
        \item Radiation testing results successfully informed critical and targeted design changes—such as the removal of problematic oscillators and MOSFETs—achieving optimal systemic reliability.
        \item With qualifications nearly complete, approximately \textbf{930 DB6 units} will enter full-scale mass production starting in 2026.
    \end{itemize}
}

\end{columns}
\end{document}
